\section{The double-ladder family}
\label{sec:double-ladders}

Let \(L_a=P_{a+1}\mathbin{\square}K_2\) be a ladder with vertices
\(A_i^r\), where \(0\le i\le a\) and \(r\in\{0,1\}\).  Its rungs are
\(A_i^0A_i^1\), and its rails are \(A_i^rA_{i+1}^r\).  Define a second,
disjoint ladder \(L_b\) on vertices \(B_j^r\).

\begin{definition}[Double ladder]
\label{def:double-ladder}
For positive integers \(a,b\), the graph \(D(a,b)\) consists of \(L_a\) and
\(L_b\) together with the four connector edges
\begin{equation}
\label{eq:connectors}
 A_0^0B_0^0,\qquad
 A_0^1B_b^0,\qquad
 A_a^0B_0^1,\qquad
 A_a^1B_b^1.
\end{equation}
Thus
\[
 |V(D(a,b))|=2(a+b)+4.
\]
\end{definition}

\subsection{The square insertion}

Suppose \(uvyxu\) is a facial quadrilateral, so that \(uv\) and \(xy\) are
opposite edges.  Replace \(uv\) by \(u p q v\), replace \(xy\) by
\(x r s y\), and add \(pr\) and \(qs\).  The four new vertices induce the
facial cycle \(p q s r p\).  In a ladder cell this operation replaces one
cell by three and increases the ladder length by two.

\begin{figure}[!htbp]
\centering
\includegraphics[width=.43\textwidth]{square_insertion_d55.png}
\caption{Annotated \(D(5,5)\).  Blue edges are rails, red edges are rungs,
orange edges are terminal attachments, and the highlighted four-vertex tile is
a certified square insertion.  The inverse operation deletes the tile and
restores the two opposite rail edges.}
\label{fig:square-insertion}
\end{figure}

\begin{lemma}[Preservation under facial square insertion]
\label{lem:square-preservation}
Let \(G\) be a simple, planar, bipartite, cubic, 3-connected graph, and apply
the preceding operation to opposite edges of a facial quadrilateral.  The
resulting graph \(G^+\) has the same five properties.
\end{lemma}

\begin{proof}
The four terminals are distinct and the inserted vertices are new, so
simplicity is immediate.  Drawing the three-cell ladder inside the old face
proves planarity.  Each old terminal loses one edge and gains one edge; every
new vertex has two neighbours on the inserted 4-cycle and one old terminal.
Thus cubicity is preserved.  If the facial quadrilateral is
\(uvyxu\), then \(u,y\) lie in one bipartition class and \(v,x\) in the other.
The unique alternating coloring along the two new length-three paths makes
both \(pr\) and \(qs\) bichromatic, proving bipartiteness.

It remains to prove 3-connectivity.  Let \(S\subseteq V(G^+)\) with
\(|S|\le2\), and call \(p,q,r,s\) the new vertices.  If \(S\) contains no new
vertex, paths in the connected graph \(G-S\) lift through the two intact
replacement paths; the intact inner 4-cycle is attached to every surviving
terminal.  If \(S\) contains exactly one new vertex, at most one old vertex is
deleted.  The graph obtained from \(G\) by deleting that old vertex is
2-connected and hence has no bridge.  Delete, if necessary, the one old edge
whose replacement path contains the removed new vertex.  The remaining old
vertices are still connected, and the other three new vertices form a path
attached to at least two surviving terminals.  Finally, if two new vertices
are deleted, no old vertex is deleted.  A 3-vertex-connected graph is
3-edge-connected, so \(G-\{uv,xy\}\) is connected.  Every remaining new
vertex has its own terminal edge into this connected old core.  Hence
\(G^+-S\) is connected in all cases, and \(G^+\) is 3-connected.
\end{proof}

\begin{theorem}
\label{thm:double-ladder-barnette}
For all odd integers \(a,b\ge1\), the graph \(D(a,b)\) is a Barnette graph.
\end{theorem}

\begin{proof}
The graph \(D(1,1)\) is isomorphic to the cube.  Increasing either odd
parameter by two is exactly the facial square insertion of
\cref{lem:square-preservation} in a cell of the corresponding ladder.  Starting
from \(D(1,1)\), repeated insertions therefore reach every pair of positive odd
parameters while preserving all Barnette properties.
\end{proof}

\subsection{Certified occurrences and face structure}

\begin{computationalresult}[Certified double-ladder recurrence]
\label{res:double-ladder-recurrence}
Exact graph isomorphisms identify the unique maximizers at orders divisible by
four as follows:
\[
\begin{array}{c|cccccccc}
n&8&12&16&20&24&28&32&36\\ \hline
G_n&D(1,1)&D(3,1)&D(3,3)&D(5,3)&D(5,5)&D(7,5)&D(7,7)&D(9,7).
\end{array}
\]
All seven consecutive transitions have reconstructible expansion certificates
whose reconstructed edge set is exactly the certified target and whose reduced
canonical hash is exactly the certified source.
\end{computationalresult}

\begin{proposition}[Face pattern in the recursive embedding]
\label{prop:double-ladder-faces}
For odd positive \(a,b\), the embedding obtained from the cube by the square
insertions above has \(a+b\) designated ladder-cell faces, all quadrilateral,
together with two faces of length \(a+3\) and two faces of length \(b+3\).
When \(a=1\) or \(b=1\), the corresponding two cap faces are themselves
quadrilateral.
\end{proposition}

\begin{proof}
For \(D(1,1)\), the two designated cells and four cap faces are the six faces
of the cube.  Increasing \(a\) by two replaces one designated cell by three,
adding two quadrilateral cells.  The two cap boundaries running along the
modified (a)-ladder each acquire two edges; the other two cap boundaries are
unchanged.  This is precisely the claimed update.  The argument for \(b\) is
identical, so induction proves the formula.
\end{proof}

For the certified members from \(D(3,3)\) onward, the two sets of designated
cells are exactly the two path components of the quadrilateral face-adjacency
graph.  At orders 8 and 12, additional quadrilateral cap faces make this
face-adjacency description degenerate, although the abstract double-ladder
definition remains valid.

\begin{figure}[!htbp]
\centering
\begin{tabular}{ccc}
\includegraphics[width=.24\textwidth]{d11.png} &
\includegraphics[width=.24\textwidth]{d53.png} &
\includegraphics[width=.24\textwidth]{d55.png} \\
\(D(1,1)\) & \(D(5,3)\) & \(D(5,5)\) \\[3pt]
\includegraphics[width=.24\textwidth]{d75.png} &
\includegraphics[width=.24\textwidth]{d77.png} &
\includegraphics[width=.24\textwidth]{d97.png} \\
\(D(7,5)\) & \(D(7,7)\) & \(D(9,7)\)
\end{tabular}
\caption{Selected certified double-ladder maximizers.  Layouts were produced
from the certified planar embeddings; geometry is used only for presentation,
while all structural identifications were made from abstract edges and rotation
systems.}
\label{fig:double-ladders}
\end{figure}
